\section{Introduction}
\label{sec:intro}

% The extension-disclosure \thanks{} footnote was REMOVED from main.tex's
% \author block 2026-07-10 (conference draft never submitted -> standalone
% first submission, no disclosure obligation; see the comment at that site).

Knowledge editing promises surgical updates to a language model's factual
associations without retraining. Locate-then-edit methods -- rank-one model
editing (ROME) \citep{meng2022rome} and its mass-editing successor MEMIT
\citep{meng2023memit} -- realize this with a rank-one or low-rank
write to a single multilayer-perceptron (MLP) down-projection at a
``factual-storage'' layer. The
write is effective at installing the target fact, but it is not local: it
also shifts the model's read-out on \emph{unrelated} facts, a phenomenon we
call collateral damage. Understanding, predicting, and removing that damage is
now the central reliability question for the method family.

\paragraph{Why editing reliability matters.} Because such updates are cheap,
they invite repeated, low-oversight application in deployed models -- and that
only holds up if each edit's collateral cost is knowable \emph{before} it is
applied. A predictor computable from quantities the editor already holds --
the edit key and the keys of the facts to preserve -- is therefore a
prerequisite for treating locate-then-edit methods as a dependable maintenance
tool rather than a benchmark curiosity.

A natural question is whether collateral damage is predictable
\emph{a priori} -- before running the edit -- from geometry alone. The
concurrent work CLaRE-ty \citep{clarety2026} answers affirmatively with an
\emph{empirical} single-layer forward-activation cosine, evaluated over
\clareNfacts{} facts, \clareNeditors{} locate-then-edit editors, and
\clareNmodels{} attention models, and -- by its authors' own statement --
correlational, with no causal mechanism. We take a different and complementary
angle: not a new predictor, but an \emph{account} of the mechanism, its
failure modes, and a causal test.

\paragraph{Scope (stated up front).} The \emph{signed} key-cosine$\rightarrow$damage
law we characterize is \emph{Llama-family-specific}. We frame the paper as a
Llama-family mechanistic account \emph{plus} a cross-architecture dissociation
arm, not a universal predictor. Two laws with different reach recur throughout:
the \emph{signed} law (sign and rank of the key-cosine coupling) holds on the
Llama family and off it either goes null (gemma, Phi) or inverts (Qwen); the
\emph{magnitude} law ($|C|\rightarrow|\text{damage}|$) transfers on four of five
non-Llama families, with gemma the documented exception.

\paragraph{Contributions.}
\begin{itemize}
\item \textbf{A closed-form \SxC{} decomposition (\S\ref{sec:method}).} ROME's
  collateral perturbation to a probe factorizes as
  $S\cdot\|k_p\|\cdot|C|$: ``cosine ranks damage'' is built into the rank-one
  algebra. The product \SxC{} is the exact closed-form reduction of
  first-order gradient influence, obtainable with no backprop -- it explains
  \emph{why} cosine works and predicts \emph{where} it stops. We give it a
  formal statement as a zero-cost \emph{rank estimator} and report its honest
  limit: it is not a faithful rank-surrogate of the true influence it reduces
  (\S\ref{sec:method}).
\item \textbf{A depth/regime transition and a cross-architecture editability
  band (\S\ref{sec:regime}).} On the confound-clean within-probe metric,
  key-cosine beats matrix-norm-growth at L8/L10/L12 and norm-growth overtakes at
  L14; the coupling's sign tracks the sign of the damage regime across scale.
  The editable depth range within which any such law can be measured is itself
  architecture-conditioned: across six base models over four families (1B--20B)
  it ranges from a wide plateau (Llama, Pythia) through a gradual decay (GPT-J)
  to a shallow-band cliff (GPT-NeoX-20B).
\item \textbf{An editor/architecture dissociation (\S\ref{sec:dissociation}).}
  The law is locate-then-edit-mechanism-specific: across six editors coupling
  and damage move oppositely, from ROME through a null-space projector to a
  codebook editor whose weight delta is identically zero; it dies for full-rank
  fine-tuning and for MEMIT, and is near-null/inverted off Llama -- each
  breakdown predicted by the same algebra. On a second architecture (GPT-J) the
  coupling reappears but relocates in depth and inverts in sign, so the law's
  depth-of-action is architecture-conditioned rather than a fixed fraction of
  network depth.
\item \textbf{A causal test and its scale/architecture generality
  (\S\ref{sec:causal}).} AlphaEdit's null-space projection removes the
  geometry-predicted damage in proportion to key-cosine at every layer, and
  erases the Qwen inversion -- the causal pathway CLaRE-ty explicitly leaves open.
  A held-out projector retires the circularity objection. The intervention
  generalizes at full strength to an instruction-tuned twin and
  sign-agnostically to a second architecture (both three-seed), and
  monotonically at 20B (single-seed),
  while the signed coupling attenuates at 8B, which we report as
  claim-tightening rather than a positive confirmation at scale.
\item \textbf{Extension to a new edit type and datasets
  (\S\ref{sec:deletion},~\S\ref{sec:generality}).} The same mechanism governs
  collateral damage from refusal-\emph{deletion} edits, and the law replicates
  across datasets -- zsRE and the multi-hop MQuAKE / MQuAKE-T benchmarks under
  an overlap-robustness audit -- an instruction-tuned twin, and a downstream
  GLUE task bridge.
\end{itemize}

\iftaslp
Framed as a language-technology maintenance problem, the contributions above
amount to a diagnostic and a remedy for a specific failure mode of automated
factual patching in deployed language models: a zero-cost pre-edit check for
which corrections are likely to be collaterally expensive, and a null-space
projection remedy that removes exactly the damage the check predicts. Neither
piece assumes access to the full post-hoc evaluation suite that auditing a
patched system would ordinarily require, which is the practical setting a
maintenance pipeline actually faces.
\fi

\paragraph{Roadmap.} Section~\ref{sec:related} positions the account;
\S\ref{sec:method} derives the decomposition and measurement protocol;
\S\ref{sec:regime}--\ref{sec:dissociation} report the depth/regime, causal,
and dissociation results; \S\ref{sec:deletion}--\ref{sec:sequential} cover
deletion edits, generality, and a sequential stress test (edit-key anisotropy
appears in the supplementary material); \S\ref{sec:discussion} discusses
implications and limitations.

Code, per-cell result artifacts, and figure-generation scripts sufficient to
reproduce every number will be released upon publication.
